\section{The index on the split threefold}\label{sec:original}

Take \(q\) as the external circle. Its fixed threefold over
\(L\in B=\operatorname{Pic}^{b}(C)\) is the total space of
\(L\oplus K_C L^{-1}\), where \(b=g-1\). The internal circle
\(T=\mathbb C^\times_t\) rotates these two lines with opposite weights.
We shall compute the class of the index on the space of cycles.
This requires keeping the dependence on \(L\). For example, a section of
\(L\) translates the curve inside its threefold, and the number of
independent translations can jump as \(L\) varies.

Write \(\mathcal L\) for the normalized Poincar\'e line on \(C\times B\),
\(\mathcal M=K_C\mathcal L^{-1}\), and
\(X=\operatorname{Tot}_{C\times B}(\mathcal L t\oplus\mathcal M t^{-1})\).
Let \(\mathcal P=P_n(X/B,m)\) parametrize stable pairs
\(\mathbb I=[\mathcal O_X\to\mathcal F]\) in the fibers of \(X\to B\).
The sheaf has proper pure one-dimensional support, generic degree \(m>0\)
over \(C\), Euler characteristic \(n\), and a section with
zero-dimensional cokernel~\cite{pt}.
Let \(h:\mathcal P\to W_m\) send a pair to its support cycle in the
fivefold, and let \(s_m:B\to W_m\) denote the cycle \(m[C]\).
The support morphism is proper. Indeed, compactifying \(X\) to the
projective bundle with infinity divisor \(D_\infty\) gives the projective
stable-pair support morphism. Its restriction over cycles disjoint
from \(D_\infty\) parametrizes pairs on \(X\), since purity forces the
whole sheaf support to lie on its cycle. The inclusion of cycles on the
fixed threefold into the fivefold cycle space is closed.

\subsection{The orientation and the support of the index}

Push the universal sheaf down to \(C\times\mathcal P\) and call the result
\(E\). The projection of the support is proper and affine, hence finite.
Purity makes every fiber of \(E\) over \(\mathcal P\) torsion-free on \(C\).
Flatness over \(\mathcal P\) then makes \(E\) a vector bundle of rank \(m\).
For the projection \(\rho:C\times\mathcal P\to\mathcal P\), set
\begin{equation}
 P=R\rho_*E-
 R\rho_*\bigl(\operatorname{End}(E)(1-\mathcal L t)\bigr).
 \label{eq:originalpolarization}
\end{equation}
This is a polarization of the relative virtual tangent class:
\begin{equation}
 T^{\rm vir}_{\mathcal P/B}=P-P^\vee,\qquad
 \operatorname{rk}P=n+m^2b .
 \label{eq:originaltangent}
\end{equation}
Indeed the two normal coordinates give a Koszul resolution for the
self-Ext term, whose class before pushforward is
\(\operatorname{End}(E)(1-\mathcal Lt)(1-\mathcal M t^{-1})\).
Duality on \(C\) pairs its two halves and pairs the two terms linear in
\(E\) in the framed deformation complex. Riemann--Roch gives the rank
in~\eqref{eq:originaltangent}. We use \((\det P)^{-1}\) as the square
root of the relative virtual canonical bundle. The two ordered external
lines are trivial as bundles in this presentation, so their determinant
ratio introduces no further line. Write
\(\widehat{\mathcal O}^{\rm vir}=(\det P)^{-1}\mathcal O^{\rm vir}\)
with the signed counting variable \(-q\)~\cite{no}.

Let \(F=\mathcal P^T\) and \(i:F\to\mathcal P\) denote its inclusion.
A fixed pair has its entire reduced support on the zero section. Indeed,
every fiber of its support over \(C\) is finite, so the connected torus
fixes each point of that fiber. The origin is the only fixed point in
each affine-plane fiber. In particular \(F\) is proper over \(B\),
by properness of the support morphism and invariance of properness
under nilpotent thickenings. Fixed supports can themselves have
nontrivial scheme structure.

To obtain an \emph{integral} lift of the virtual sheaf to \(F\), consider
the vector field induced by the normal rotation. Relative Calabi--Yau
duality identifies \(\operatorname{Ob}_{\mathcal P/B}\) with
\(\Omega_{\mathcal P/B}\), and contraction with the vector field gives
a cosection. We check that it descends to the absolute obstruction
sheaf. Write \(\xi\) for the Euler field on \(X\) and
\(D=\iota_\xi\operatorname{At}_{X/B}(\mathbb I)\) for the corresponding
degree-one endomorphism. A Picard tangent represented by
\(\alpha\in H^1(C,\mathcal O_C)\) varies the transition functions of
\(\mathcal L\) and \(\mathcal M\) inversely. Its Kodaira--Spencer class
is therefore \(\alpha\xi\), and its obstruction is the trace-free part
of \(\alpha D\), up to the common infinitesimal-action sign~\cite{ht}.
Its pairing with \(D\) vanishes. In the Serre trace, commuting \(D\)
past the even block \(\alpha D\) preserves the sign, whereas commuting
it past the odd scalar \(\alpha\) changes the sign. This calculation
also applies over nonreduced parameter schemes. For the proper Serre
trace one can use the projective bundle compactification. The trace-free
endomorphism complex vanishes near infinity, where \(\mathbb I\) is the
identity line. The relative canonical bundle on the compactification is
\(\mathcal O(-3D_\infty)\), and the section of
\(\mathcal O(3D_\infty)\) identifies the twisted trace-free complex
with the untwisted one in the derived category. Multiplication by this
section has degree zero and leaves the cyclic sign unchanged.

Thus the cosection kills the image of \(T_B\) in the relative obstruction
sheaf and descends through the absolute obstruction sequence. Its zero
scheme is \(F\): on a homogeneous function of weight \(j\) the orbit
derivation is multiplication by \(j\), and the nonzero weights
generate the fixed-scheme ideal in characteristic zero.
Equivariant cosection localization now gives
\begin{equation}
 \gamma\in G_0^T(F),\qquad i_*\gamma=[\mathcal O_{\mathcal P}^{\rm vir}].
 \label{eq:originalintegrallift}
\end{equation}
Here \(G_0^T\) denotes the group of equivariant coherent sheaves with
integer coefficients. We use the absolute virtual sheaf induced by the
relative theory over the smooth base and the relative orientation
fixed above~\cite{kiemli,kiemsavvas}.

\subsection{The fixed class over the Picard variety}

On \(F\), decompose \(P=P_0+P_{\rm mov}\) into its zero and nonzero
weights. The moving virtual normal class is
\(P_{\rm mov}-P_{\rm mov}^\vee\).
To compute \(\gamma\) without losing torsion in the Picard family,
choose a two-term bundle resolution
\([N^{-1}\to N^0]\) of the moving cotangent complex on \(F\).
One can choose it weight by weight since \(F\) is quasi-projective and
\(T\) acts trivially on it. For \(e(V)=\lambda_{-1}(V)\), the
self-intersection and composition identities in virtual pullback
give~\cite[proof of Theorem 3.3 and Lemma 3.6]{qu}
\begin{equation}
             e(N^0)\gamma=e(N^{-1})[\mathcal O_F^{\rm vir}].
 \label{eq:originalintegralcancel}
\end{equation}
Indeed virtual self-intersection applied to \(i_*\gamma\) gives the
left side. Applied to the virtual sheaf on \(\mathcal P\), the same
pullback adjoins \(N^{-1}[1]\) to the fixed obstruction theory and
gives the right side. Both operations precede division by torus
characters: the factors in~\eqref{eq:originalintegralcancel} are the
Koszul classes for zero sections.

On each component a moving bundle has Euler polynomial in \(t\) with
invertible extremal coefficients, up to sign and a line bundle.
Hence multiplication by it on \(G_0(F)[t^{\pm1}]\) is injective
even in the presence of additive torsion. The identity
\(1-a=-a(1-a^{-1})\) applied to the lines of a splitting of the
moving polarization gives
\[
 \frac{e(N^{-1})}{e(N^0)}
     =(-1)^{\operatorname{rk}P_{\rm mov}}\det P_{\rm mov}.
\]
Clearing the Euler factors and using injectivity in
\eqref{eq:originalintegralcancel} identifies the oriented lift
integrally as
\begin{equation}
 (\det P)^{-1}\gamma=
 (-1)^{\operatorname{rk}P_{\rm mov}}
       (\det P_0)^{-1}[\mathcal O_F^{\rm vir}].
 \label{eq:originalfixedlift}
\end{equation}
In particular its internal rotation weight is zero.

It remains to determine how this class varies with \(L\).
Fix \(L_0\in B\) and put \(A=\mathcal L\otimes L_0^{-1}\).
An equivariant stable pair on the fiber at \(L_0\) gives a pair on the
varying threefold by tensoring each weight-\(j\) summand \(E_j\) with
\(A^j\). Both normal actions have precisely the required changes of
weight and the section has weight zero. The inverse construction uses
\(A^{-1}\). The degree of \(A\) on each curve is zero, so this preserves
\(n\). Stable pairs have trivial automorphisms and a fixed family has a
unique compatible linearization. The construction commutes with
arbitrary base change and gives \(F\simeq F(L_0)\times B\).

This identification includes the fixed obstruction theory. In local
frames of \(A\), the weight-zero relative derived endomorphisms, the
trace and the maps induced by the section glue with total power
\(A^0\). The same identification respects their obstruction morphism:
the relative Atiyah class in the pair-parameter directions is taken
with \(A\) held fixed on \(C\times B\).
Thus the relative fixed virtual sheaf is pulled back from \(F(L_0)\).
The terms of \(P_0\) also have total power \(A^0\), and the moving
rank in~\eqref{eq:originalfixedlift} is constant on each product
component. Proper pushforward of the oriented fixed class therefore
gives a numerical coefficient \(c_{m,n}\) times \([\mathcal O_B]\).
On \(F_{\rm red}\) the support map factors through \(s_m\). Filtering
coherent sheaves by the nilradical permits the same pushforward
calculation even if the map on \(F\) has infinitesimal components
normal to \(s_m(B)\). Consequently
\begin{equation}
  h_*[\widehat{\mathcal O}_{\mathcal P}^{\rm vir}]
     =(s_m)_*\bigl(c_{m,n}[\mathcal O_B]\bigr)
     \quad\text{in }G_0^T(W_m).
 \label{eq:originalsupport}
\end{equation}

It suffices now to evaluate \(c_{m,n}\) on one fiber. The antidiagonal
local-curve stable-pair formula at levels \((b,b)\)
gives~\cite{monavari,refined}
\begin{equation}
 \sum_n c_{m,n}(-q)^n
   =\sum_{\lambda\vdash m}\ \prod_{u\in\lambda}
     \bigl((1-q^{h(u)})(1-q^{-h(u)})\bigr)^b ,
 \label{eq:originalhook}
\end{equation}
where \(h(u)\) is one plus the number of boxes to the right and below
\(u\) in its row and column. The sign follows from
\((1-q^h)(1-q^{-h})=-(q^{h/2}-q^{-h/2})^2\). In degree one this
gives \(((1-q)(1-q^{-1}))^b\). For \(g\geq1\) every summand is an
integral Laurent polynomial. Equation~\eqref{eq:originalsupport}
identifies the family class to which this numerical formula belongs.
